% \documentclass[draft]{report}

\RequirePackage{luatex85}% TeXLive 2017 fix for \geometry
\documentclass{kdp}
\usepackage{geometry}

% \usepackage{algorithm}% http://ctan.org/pkg/algorithms
\usepackage{algpseudocode}% http://ctan.org/pkg/algorithmicx
% \usepackage{algorithm}% http://ctan.org/pkg/algorithms
% \usepackage{algpseudocode}
\usepackage{tlatex}
\usepackage{listings}
\usepackage{xcolor}
\usepackage{comment}
\usepackage{fancyhdr}
\usepackage{amssymb}
\usepackage{inputenc}
\usepackage{svg}
\usepackage{framed}
\usepackage{graphicx}
\usepackage{tikz}
\usepackage{minted}
\setminted{fontsize=\small} % applies to all minted environments
\usetikzlibrary{automata, positioning, arrows}

\usepackage{hyperref}
\hypersetup{
    colorlinks,
    citecolor=black,
    filecolor=black,
    linkcolor=black,
    urlcolor=black
}

\pagestyle{fancy}
\fancyhf{}

\fancyhead[L]{\nouppercase{\leftmark}} % Chapter number and title on the left
\fancyhead[R]{\thepage}     % Right-aligned header

\lstset { %
    language=C++,
    backgroundcolor=\color{black!5}, % set backgroundcolor
    basicstyle=\footnotesize,% basic font setting
}

\title{Scalable Software with Haskell \\ Early Preview}

% \maketitle

\author{Richard Tang}
\date{\today}
\begin{document}
\maketitle

\section*{Acknowledgement}

\tableofcontents

% # Functional language Haskell/Elixir

% [Optimizing Haskell Programs](https://sriramsami.com/haskell-optimization/)

% Other compiled memory safe language?

% Facebook spam filter implementation? Reference to hackernews of how it integrates with cpp

% Haskell garbage collector?

% Functor, applicative, monad

% Packet processing

% Multicore parralellization

% No datashare

% Representing graph, tree, etc in Haskell

% Integrate with c or cpp

% Disassembly with readlf and objdump

% High variance language

% https://news.ycombinator.com/item?id=19311750

% https://clouddevs.com/elixir/high-performance-computing/

% Can elixir be formally verified being a functional language?

% Byte code translation provides fault tolerance but has runtime overhead

% [A Million-user Comet Application with Mochiweb, Part 3 | Richard Jones](https://www.notion.so/A-Million-user-Comet-Application-with-Mochiweb-Part-3-Richard-Jones-1ab572a03bb38150af49f6002d44fe2b?pvs=21)

\chapter{Introduction}

Functional language is declarative in nature and by definition has no side
effect. Imperative language such as C permits mutating the global state, which
increases comprehension complexity. Consider the following example:

\begin{minted}{c}

    int eventLoop() { 
        int evt = receiveEvent();
        switch (evt) { 
            case EVENT_A: {
                processEventA(evt);
                break;
            }
            case EVENT_B:
            case EVENT_C:
        }
    }

    int processEventA(int evt) { 
        /* processing */ 
        reportEvent(evt);
        global.state = PROCESSED_A; 
    }
\end{minted}

Changing the \textit{global.state} may alter the system in scope larger than
the function itself. \textit{global} can be either file or program scope, but
either way publicly accessible to other parts of the program. Other parts of
the program may alter either local or global behavior based on this variable.
For example, eventLoop can condition on global.state and stops processing
event. Imperative language like C does not provide any safe guard against this,
as it simply provides an abstraction of the hardware, which is also imperative
in nature.\\

This increases the burden to reason about program correctness, and allows the
program more susceptible to mistakes. C/C++ heavily encourages the use of
\textit{const}, which is meant to lighten the cost of these issues.
Fundamentally, C/C++ is designed to offer low-level control, so \textit{const}
is a best practice but not the must. On a related note, Rust defaults all
variables to be immutable and variables need to be declared with \textit{mut}
to mutate. In any case, these reduces the likelihood of the problems described
above, but does not eliminate them completely. The language fundamentally does
not guarantee functional pureness, as it not an design objective.

\section{Composability}

Functional languages, on the other hand, promises \textit{no side-effects}. For
a set of inputs, a function always produce the same output. This also allows
\textit{composability}, akin to chaining UNIX utilities with pipe:

\begin{minted}{bash}
    cat input.txt | grep -rn needle | wc
\end{minted}

\begin{minted}{haskell}
    let rv = count . (filter "needle") lines 
\end{minted}

Since the definitions are functional, Haskell compiler may collapse chained
functions underneath the hood to provide an optimized compiled
implementation.\\

\section{Horizontal Scalability}

Another attractive property of functional programming is horizontal
scalability. By eliminating shared variables entirely, the system becomes
\textit{much} easier to reason about parallel behavior. While locks such as
mutex are convenient to ensure atomicity during shared variable access, they
introduce secondary order problem such as performance reasoning. Any shared
resources are subjected to contention, which can be a source a bugs or subtle
throughput bottlenecks.\\

TODO: insert anecdote on throughput debug via over intensive locking.

\section{Memory Safe Language}

\section{Multicore Safety}

\section{Developer Efficiency}

\section{Expressiveness}

To provide some context around the language expressiveness, we will walk
through a simple socket programming example. Consider a socket programming
example with the following requirements:

\begin{itemize}
    \item Create a socket
    \item Serialize the payload
    \item Send the payload
    \item Receive the payload
    \item Deserialize the payload
    \item Process any error handling
\end{itemize}

The above requirement can be described by the following code:

\begin{minted}{haskell}
dispatch :: Integer -> DRequest -> Chan (RequestResponse) -> IO ()
dispatch port req sq = do
  bracket
    (socket AF_INET Stream defaultProtocol)
    (\sock -> close sock)
    (\sock -> do
       connect sock (SockAddrInet (fromIntegral port) 0)
       resp <-
         handle (\(_ :: SomeException) -> return $ CRespError) $ do
           sendAll sock (DBL.toStrict $ serialize req)
           maybeResp <-
             runMaybeT
               $ liftIO (recv sock 4096)
                   >>= MaybeT . return . deserialize . DBL.fromStrict
           return $ fromMaybe (CRespError) (maybeResp)
       writeChan sq $ CResponse resp)
\end{minted}

The function uses two patterns, \textit{bracket} and \textit{handle}.
\textit{bracket} provides a template to insert init, body and de-init
functions. In this example, the init function creates a socket, and the de-init
function closes the socket. This approximates RAII in C++. Inside the body, we
use the \textit{handle} pattern, which returns \textit{CRespError} on any
exception (eg. from either send or receive).\\

Receiving and parsing the payload is a bit more complex, and will be described
in detail below.\\

\textit{recv} returns an \textit{IO ByteString} and is lifted into
\textit{MaybeT IO ByteString} with liftIO. liftIO is contextual; in this case
the type inference system lifts it into \textit{MaybeT IO ByteString} because
the type on the other side of the bind operator.\\

\textit{MaybeT . return . deserialize . DB.fromStrict} is a composited function.
The bind operator extracts the value from the context and returns a ByteString.
The \textit{ByteString} then goes through a series of transformation:
\begin{itemize}
    \item \textit{fromStrict} transforms ByteString to strict ByteString
    \item \textit{deserialize} transforms strict ByteString to Maybe A
    \item \textit{return} transforms Maybe A to IO Maybe A
    \item \textit{MaybeT} transforms IO Maybe A to MaybeT IO A
\end{itemize}

Finally, \textit{runMaybeT} trasnforms \textit{MaybeT IO A} to \textit{IO Maybe
A}, and the bind arrow operator unwraps and assigns \textit{Maybe A} to 
\textit{maybeResp}.

\part{Fundamental}

\chapter{Breath-First Search}

\include{bfs}

\chapter{Depth-First Search}

\include{dfs}

\chapter{Union-Find}

\chapter{Binary Search}

\chapter{Recursion}

\chapter{Tree}

\chapter{Graph}

\chapter{Functor}

\part{More Complicated}

\chapter{Applicative}

\chapter{Monad}

\chapter{System Call}

\part{Underneath the hood}

\chapter{Disassembly}

\chapter{Memory Footprint}

\chapter{Storage Footprint}

\chapter{Runtime Threads}

\chapter{Garbage Collector}

\chapter{Vectorization}

\chapter{Foreign Function Interface}

\part{Scalable Software}

\chapter{Reference}

\end{document}
